\chapter{Standard Denavit--Hartenberg Parameters}
\label{ch:dh}

\section{Convention and geometry}
The project derives its table using \textbf{standard DH}:
\begin{equation}
 A_i=R_z(\theta_i)\,T_z(d_i)\,T_x(a_i)\,R_x(\alpha_i),
 \qquad {}^{i-1}T_i=A_i.
 \label{eq:standard-dh-order}
\end{equation}
The revolute coordinate $\theta_i$ rotates about $z_{i-1}$. After this
rotation, the offset $d_i$ moves along that axis. The common-normal displacement
$a_i$ is along $x_i$; finally $\alpha_i$ rotates about that common-normal axis
to align the next joint axis. Because the transforms act right-to-left on a
point, the matrix product is written in the order in \cref{eq:standard-dh-order}.
This is not modified DH; swapping $T_x$ and $R_x$ or using the modified-DH
table would change the kinematics.

For parameters $(\theta,d,a,\alpha)$, the homogeneous transform is
\begin{equation}
 A(\theta,d,a,\alpha)=
 \begin{bmatrix}
 c_\theta&-s_\theta c_\alpha&s_\theta s_\alpha&a c_\theta\\
 s_\theta&c_\theta c_\alpha&-c_\theta s_\alpha&a s_\theta\\
 0&s_\alpha&c_\alpha&d\\0&0&0&1
 \end{bmatrix}.
 \label{eq:dh-general}
\end{equation}
The $\alpha$ and $d$ values are signed coordinate displacements and rotations;
they need not be positive physical link lengths. Frame axes must form
right-handed orthonormal triads, with $y_i=z_i\times x_i$.

\section{Mapping URDF axes to DH axes}
All six arm joint origins in the simulation URDF have zero RPY. Their joint
origin positions and local axes are in \cref{tab:arm-joints}. At the zero
configuration, the DH joint axes are chosen as
\begin{align}
 z_0&=+Z, &z_1&=+Y, &z_2&=-Y,\\
 z_3&=-X, &z_4&=-Y, &z_5&=-X.
\end{align}
These choices reproduce the positive directions in the URDF. The corresponding
local joint axes and URDF origins are listed in \cref{tab:arm-joints}. DH frames are
free to differ from link frames as long as each $A_i$ maps correctly; the
table's $q_i$ remain the URDF joint positions in the same order.

\section{Frame-by-frame construction}
\begin{enumerate}[label=\textbf{Joint \arabic*:},leftmargin=*]
 \item \textbf{Base rotation S.} $z_0=+Z$ is the base yaw axis. $z_1=+Y$ is
 the shoulder axis at $(0.040,0,0.330)$ m. Their axes are perpendicular and
 offset by the base/shoulder geometry. Choose $x_1=+X$, the common normal
 direction. The vertical displacement is $d_1=0.330$ m and the common-normal
 offset is $a_1=0.040$ m. The $-\pi/2$ twist aligns the next z-axis with
 $+Y$. The joint angle is $\theta_1=q_1$.

 \item \textbf{Shoulder L.} $z_1=+Y$ and $z_2=-Y$ are antiparallel axes.
 The upper-arm separation is $0.445$ m along home $+Z$, so choose
 $x_2=+Z$ and $a_2=0.445$ m. The antiparallel axis choice gives
 $\alpha_2=\pi$. The URDF zero shoulder angle must place $x_2$ along $+Z$;
 with the chosen preceding frame this requires the fixed offset
 $\theta_2=q_2-\pi/2$. There is no $d_2$ translation.

 \item \textbf{Elbow U.} $z_2=-Y$ and $z_3=-X$ define perpendicular axes.
 The short elbow offset is $0.040$ m along $+Z$ (the common normal), giving
 $a_3=0.040$ m. The chosen twist is $\alpha_3=-\pi/2$ and the angle is
 $\theta_3=q_3$. There is no z-offset.

 \item \textbf{Wrist roll R.} $z_3=-X$ and $z_4=-Y$ intersect at the wrist
 axis geometry after moving to the wrist center. The center is displaced
 $+0.440X$ from frame 3, opposite $z_3$; hence the signed parameter is
 $d_4=-0.440$ m. Since the axes intersect, $a_4=0$. The table uses
 $\alpha_4=+\pi/2$ and $\theta_4=q_4$.

 \item \textbf{Wrist bend B.} $z_4=-Y$ and $z_5=-X$ intersect at the same
 wrist center, so $a_5=d_5=0$. The selected twist is $\alpha_5=-\pi/2$ and
 the revolute coordinate is $\theta_5=q_5$.

 \item \textbf{Wrist twist T.} $z_5$ and $z_6$ are both $-X$ at the same
 point. There is no translation or twist, so $a_6=d_6=\alpha_6=0$, with
 $\theta_6=q_6$.
\end{enumerate}

The origin/axis picture is in \cref{fig:dh-frames}. It is a home-pose X--Z
projection; axes parallel to base $Y$ are marked out of or into the page.
The signed offset $d_4=-0.440$ is particularly easy to misread: the actual
hardware separation is positive length, but the selected $z_3$ coordinate
points opposite the displacement.

\section{Verified DH table}
\begin{table}[htbp]
\centering
\caption{Verified standard-DH parameters for the task URDF. Angles are in
radians, distances in metres.}
\label{tab:dh}
\begin{tabular}{c l r r r l}
\toprule $i$ & axis & $a_i$ & $\alpha_i$ & $d_i$ & $\theta_i$\\\midrule
1&S&0.040&$-\pi/2$&0.330&$q_1$\\
2&L&0.445&$\pi$&0&$q_2-\pi/2$\\
3&U&0.040&$-\pi/2$&0&$q_3$\\
4&R&0&$\pi/2$&-0.440&$q_4$\\
5&B&0&$-\pi/2$&0&$q_5$\\
6&T&0&0&0&$q_6$\\\bottomrule
\end{tabular}
\end{table}

Only the geometry columns $a_i,\alpha_i,d_i$ are constant. Each $q_i$ is
variable, and the sole constant angular correction is $-\pi/2$ in the second
row. The mapping preserves the URDF joint names and signs: there are no
additional sign inversions.

\section{Derivation of all six matrices}
Define $c_i=\cos q_i$, $s_i=\sin q_i$, and
$v_2=q_2-\pi/2$. Substitution into \cref{eq:dh-general} yields
\begin{align}
A_1&=\begin{bmatrix}
c_1&0&-s_1&0.040c_1\\s_1&0&c_1&0.040s_1\\0&-1&0&0.330\\0&0&0&1
\end{bmatrix},\\[1mm]
A_2&=\begin{bmatrix}
\cos v_2&\sin v_2&0&0.445\cos v_2\\
\sin v_2&-\cos v_2&0&0.445\sin v_2\\
0&0&-1&0\\0&0&0&1
\end{bmatrix},\\[1mm]
A_3&=\begin{bmatrix}
c_3&0&-s_3&0.040c_3\\s_3&0&c_3&0.040s_3\\0&-1&0&0\\0&0&0&1
\end{bmatrix},\\[1mm]
A_4&=\begin{bmatrix}
c_4&0&s_4&0\\s_4&0&-c_4&0\\0&1&0&-0.440\\0&0&0&1
\end{bmatrix},\\[1mm]
A_5&=\begin{bmatrix}
c_5&0&-s_5&0\\s_5&0&c_5&0\\0&-1&0&0\\0&0&0&1
\end{bmatrix},\\[1mm]
A_6&=\begin{bmatrix}
c_6&-s_6&0&0\\s_6&c_6&0&0\\0&0&1&0\\0&0&0&1
\end{bmatrix}.
\label{eq:all-dh-matrices}
\end{align}
For example, row 1 has $\alpha_1=-\pi/2$, so
$\cos\alpha_1=0$ and $\sin\alpha_1=-1$. Inserting those into the general
matrix gives its second and third columns and confirms the negative sign in
the $(3,2)$ entry. Row 2 uses $\cos\pi=-1$, which explains its negative
$\cos v_2$ entry. This substitution process is the safest way to catch common
DH sign errors.

\section{Base and tool transforms}
The world-to-base and base-to-DH-frame-0 transforms are identities in this
spawn configuration. Link \code{link_6_t} to the selected DH frame 6 has
rotation
\begin{equation}
 C=\begin{bmatrix}0&0&-1&0\\0&1&0&0\\1&0&0&0\\0&0&0&1\end{bmatrix}.
\end{equation}
The fixed transform from frame 6 to grasp frame is constructed from the exact
URDF fixed joints:
\begin{equation}
 {}^6T_g=C^{-1}T_{link6,flange}T_{flange,tool0}T_{tool0,g}.
 \label{eq:dh-tool-transform}
\end{equation}
Numerically, using rounded values for display,
\begin{equation}
 {}^6T_g\approx\begin{bmatrix}
 1&0&0&-0.0009111\\0&-1&0&-0.0123\\0&0&-1&-0.22598\\0&0&0&1
 \end{bmatrix}.
\end{equation}
Then the chain actually compared against runtime FK is
\begin{equation}
 {}^BT_g(q)=A_1A_2A_3A_4A_5A_6\,{}^6T_g.
 \label{eq:dh-chain}
\end{equation}

\section{Independent derivation versus executable code}
The runtime \code{Kinematics.fk} does \emph{not} evaluate the six hard-coded
DH matrices. It traverses the URDF chain and multiplies each joint origin by
the rotation about that joint's normalized axis. The standalone
\file{docs/verify_dh.py} separately defines the table and matrix in
\cref{eq:dh-general}, builds the fixed tool transform, and compares its result
to URDF FK. This separation is valuable: if both computations reused the same
table implementation, agreement would be a weaker check.

The verifier also checks the URDF joint origins, zero RPY, and local axes
against expected source geometry. The DH table was chosen to fit this exact
task URDF. It has not been proven to reproduce the manufacturer CAD or the
other URDF exports stored under the SolidWorks directory.
